% ================================================================
%  conteudo/equipamentos.tex
%  Descrição detalhada dos equipamentos do CIA/FURG:
%  histórico, princípio, cálculos e importância.
% ================================================================

\section{Equipamentos do CIA/FURG}

O Centro Integrado de Análises da FURG (CIA/FURG) reúne infraestrutura
analítica de alto nível, oferecendo suporte a pesquisas nas áreas de
química, engenharia, ciências dos materiais e biotecnologia. A seguir
são descritos os principais equipamentos disponíveis, seus fundamentos
teóricos, metodologias de cálculo e relevância científico-industrial.

% ---------------------------------------------------------------
\subsection{GC-MS --- Cromatografia Gasosa acoplada à Espectrometria de Massas}
% ---------------------------------------------------------------

\subsubsection*{Histórico}

A cromatografia gasosa foi introduzida por Archer Martin e Richard Synge
no início da década de 1950, trabalho que lhes rendeu o Prêmio Nobel de
Química em 1952. O acoplamento com a espectrometria de massas tornou-se
viável comercialmente na década de 1960, quando McLafferty e Holmes
demonstraram que os fragmentos iônicos gerados no espectrômetro poderiam
ser usados como impressão digital molecular. Desde então, o GC-MS
consolidou-se como a técnica de referência para identificação e
quantificação de compostos orgânicos voláteis e semivoláteis em matrizes
complexas.

\subsubsection*{Princípio de funcionamento}

A amostra é vaporizada e introduzida em um fluxo de gás de arraste inerte
(geralmente He ou $\text{N}_2$) que a conduz por uma coluna capilar
recoberta internamente com uma fase estacionária. Cada componente da
mistura interage de forma distinta com essa fase, sendo eluído em tempos
diferentes --- o \textbf{tempo de retenção} ($t_R$). Ao sair da coluna, os
compostos entram no espectrômetro de massas, onde são ionizados
(tipicamente por ionização por elétrons, EI, a 70 eV) e acelerados em
direção ao analisador (quadrupolo, armadilha iônica ou TOF). O detector
registra a razão massa/carga ($m/z$) dos fragmentos, gerando o espectro de
massas, que é comparado a bibliotecas como a NIST para identificação.

\subsubsection*{Grandezas e cálculos principais}

\textbf{Fator de retenção ($k$):}
\begin{equation}
    k = \frac{t_R - t_M}{t_M}
\end{equation}
onde $t_R$ é o tempo de retenção do analito e $t_M$ é o tempo morto da
coluna (tempo de retenção de um composto não retido).

\textbf{Resolução entre dois picos ($R_s$):}
\begin{equation}
    R_s = \frac{2\,(t_{R2} - t_{R1})}{w_1 + w_2}
\end{equation}
onde $w_1$ e $w_2$ são as larguras de base dos picos. Valores $R_s \geq 1{,}5$
indicam separação completa (linha de base).

\textbf{Quantificação por padronização interna:}
\begin{equation}
    C_{\text{analito}} = \frac{A_{\text{analito}}}{A_{\text{padrão}}}
                          \times \frac{C_{\text{padrão}}}{FR}
\end{equation}
onde $A$ representa a área do pico cromatográfico e $FR$ é o fator de
resposta relativo do detector para o analito em relação ao padrão interno.

\subsubsection*{Importância}

O GC-MS é fundamental na análise de ácidos graxos (como ésteres metílicos
de ácidos graxos, FAMEs), compostos de aroma, contaminantes ambientais,
fármacos e metabólitos de fermentação. No contexto do CIA/FURG, a técnica
é especialmente relevante para caracterização de biomassa microalgal e
bioprodutos, diretamente vinculada ao perfil de pesquisa da engenharia
bioquímica da instituição.

% ---------------------------------------------------------------
\subsection{BET --- Adsorção de Nitrogênio (Área Superficial)}
% ---------------------------------------------------------------

\subsubsection*{Histórico}

O modelo BET foi desenvolvido em 1938 por Stephen Brunauer, Paul Emmett
e Edward Teller como extensão da teoria de Langmuir para adsorção em
multicamadas. Publicado no \textit{Journal of the American Chemical Society},
tornou-se o método padrão internacional (ISO 9277) para determinação de
área superficial específica de sólidos porosos.

\subsubsection*{Princípio de funcionamento}

A amostra é submetida a vácuo e resfriada a 77 K (temperatura do
nitrogênio líquido). Doses controladas de $\text{N}_2$ são introduzidas
progressivamente, e a quantidade adsorvida é medida manométrica ou
gravimetricamente em função da pressão relativa ($P/P_0$).

\subsubsection*{Grandezas e cálculos principais}

A equação linearizada BET é:
\begin{equation}
    \frac{P/P_0}{n\,(1 - P/P_0)}
    = \frac{1}{n_m\,C} + \frac{C-1}{n_m\,C}\cdot\frac{P}{P_0}
\end{equation}
onde $n$ é a quantidade adsorvida, $n_m$ é a capacidade de monocamada e
$C$ é a constante BET relacionada à entalpia de adsorção.

A área superficial específica é então:
\begin{equation}
    S_{\text{BET}} = \frac{n_m \cdot N_A \cdot \sigma}{m}
\end{equation}
sendo $N_A$ o número de Avogadro, $\sigma = 0{,}162\;\text{nm}^2$ a área
ocupada por uma molécula de $\text{N}_2$ em monocamada e $m$ a massa da
amostra.

\subsubsection*{Importância}

A área superficial é parâmetro crítico em catálise heterogênea, adsorção
de poluentes, caracterização de carvões ativados, zeólitas e materiais
nanoporosos. Em biotecnologia, é usada para avaliar suportes de
imobilização enzimática e biomateriais para liberação controlada de
fármacos.

% ---------------------------------------------------------------
\subsection{RMN --- Ressonância Magnética Nuclear}
% ---------------------------------------------------------------

\subsubsection*{Histórico}

O fenômeno da ressonância magnética nuclear foi descoberto independentemente
por Felix Bloch e Edward Purcell em 1946, ambos premiados com o Nobel de
Física em 1952. O desenvolvimento dos espectrômetros de transformada de
Fourier (FT-RMN) por Richard Ernst na década de 1970 (Nobel de Química,
1991) revolucionou a técnica, tornando-a a principal ferramenta de
elucidação estrutural de moléculas orgânicas.

\subsubsection*{Princípio de funcionamento}

Núcleos atômicos com spin não nulo (como $^1$H e $^{13}$C) possuem momento
magnético. Quando expostos a um campo magnético estático $B_0$, alinham-se
paralelamente ou antiparalelamente a ele, com diferença de energia:
\begin{equation}
    \Delta E = h\,\nu_0 = \hbar\,\gamma\,B_0
\end{equation}
onde $\gamma$ é a razão giromagnética do núcleo. Pulsos de radiofrequência
excitam os núcleos, e o sinal de relaxação (FID --- \textit{free induction
decay}) é coletado e convertido em espectro por transformada de Fourier.

\subsubsection*{Grandezas e cálculos principais}

O \textbf{deslocamento químico} ($\delta$, em ppm) normaliza a frequência
de ressonância em relação a um padrão (TMS):
\begin{equation}
    \delta = \frac{\nu_{\text{amostra}} - \nu_{\text{ref}}}{\nu_{\text{ref}}} \times 10^6
\end{equation}

A \textbf{constante de acoplamento} $J$ (em Hz) descreve a interação
spin--spin entre núcleos vizinhos e é independente do campo $B_0$,
auxiliando na determinação da conectividade molecular.

\subsubsection*{Importância}

A RMN permite elucidação completa de estruturas orgânicas, análise de
pureza, estudo de dinâmica molecular e quantificação sem calibração
prévia (RMN quantitativa, qRMN). É indispensável em síntese orgânica,
desenvolvimento de fármacos e caracterização de polímeros e materiais.

% ---------------------------------------------------------------
\subsection{DSC --- Calorimetria Exploratória de Varredura}
% ---------------------------------------------------------------

\subsubsection*{Histórico}

A calorimetria diferencial de varredura foi introduzida comercialmente em
1963 pela Perkin-Elmer, a partir de trabalhos de Watson e O'Neill. O
método evoluiu da análise térmica diferencial (DTA), existente desde o
século XIX, mas passou a medir diretamente o fluxo de calor, e não apenas
diferenças de temperatura.

\subsubsection*{Princípio de funcionamento}

A amostra e uma referência inerte são aquecidas (ou resfriadas) a uma taxa
constante $\beta = \mathrm{d}T/\mathrm{d}t$. O instrumento mede a
diferença de fluxo de calor necessária para manter ambas na mesma
temperatura. Transições de fase (fusão, cristalização, desnaturação) e
reações químicas (polimerização, decomposição) produzem picos endotérmicos
ou exotérmicos no termograma.

\subsubsection*{Grandezas e cálculos principais}

A entalpia de transição é obtida pela integração do pico:
\begin{equation}
    \Delta H = \frac{1}{m} \int_{T_1}^{T_2}
               \frac{\dot{q}_{\text{amostra}} - \dot{q}_{\text{ref}}}{\beta}\;\mathrm{d}T
\end{equation}
A capacidade calorífica da amostra é:
\begin{equation}
    c_p = \frac{\dot{q}}{m\,\beta}
\end{equation}

\subsubsection*{Importância}

O DSC é fundamental para caracterização de polímeros (grau de
cristalinidade, temperatura de transição vítrea $T_g$), estabilidade de
fármacos, análise de lipídios e estudo de proteínas. Em engenharia de
processos, auxilia no dimensionamento de trocadores de calor e na avaliação
de riscos de reações exotérmicas.

% ---------------------------------------------------------------
\subsection{TGA --- Termogravimetria}
% ---------------------------------------------------------------

\subsubsection*{Histórico}

A termogravimetria tem raízes nos estudos de Le Chatelier no final do
século XIX, mas ganhou formato analítico consolidado na década de 1950,
com a popularização de balanças termogravimétricas de alta precisão. A
norma ASTM E1131 regulamenta atualmente o método.

\subsubsection*{Princípio de funcionamento}

A massa da amostra é monitorada continuamente enquanto a temperatura é
variada segundo um programa controlado (aquecimento dinâmico ou
isotérmico), geralmente sob atmosfera controlada (ar, $\text{N}_2$ ou
$\text{O}_2$). Perdas de massa são associadas a eventos como desidratação,
decomposição e oxidação.

\subsubsection*{Grandezas e cálculos principais}

A perda de massa percentual em cada evento:
\begin{equation}
    \Delta m\,(\%) = \frac{m_i - m_f}{m_i} \times 100
\end{equation}
A curva derivada (DTG) localiza com precisão a temperatura de máxima taxa
de decomposição:
\begin{equation}
    \text{DTG} = \frac{\mathrm{d}m}{\mathrm{d}T}
\end{equation}

\subsubsection*{Importância}

TGA é amplamente usada para determinar estabilidade térmica de polímeros,
teor de cinzas, umidade residual, composição de compósitos e cinética de
decomposição. Em biotecnologia, é aplicada à caracterização de biomassa,
avaliação de biocombustíveis e estudo de materiais encapsulantes.

% ---------------------------------------------------------------
\subsection{FTIR-Raman --- Espectroscopia Vibracional}
% ---------------------------------------------------------------

\subsubsection*{Histórico}

A espectroscopia no infravermelho por transformada de Fourier (FTIR)
popularizou-se nos anos 1970, substituindo os espectrômetros dispersivos
de prisma. O efeito Raman, por sua vez, foi previsto teoricamente por
Smekal em 1923 e observado experimentalmente por C.V. Raman em 1928
(Nobel de Física, 1930). O acoplamento de ambas as técnicas em um único
instrumento tornou-se padrão laboratorial a partir dos anos 1990.

\subsubsection*{Princípio de funcionamento}

\textbf{FTIR:} a amostra absorve radiação infravermelha em frequências que
coincidem com modos vibracionais de ligações químicas (estiramento e
deformação). O interferograma resultante é convertido em espectro por
transformada de Fourier.

\textbf{Raman:} luz monocromática (laser) ilumina a amostra. A maior parte
é espalhada elasticamente (Rayleigh), mas uma fração sofre espalhamento
inelástico com deslocamento de frequência proporcional à energia dos modos
vibracionais da molécula (deslocamento Raman, $\Delta\tilde{\nu}$, em
cm$^{-1}$).

\subsubsection*{Grandezas e cálculos principais}

O número de onda de absorção/espalhamento:
\begin{equation}
    \tilde{\nu} = \frac{1}{\lambda} \quad [\text{cm}^{-1}]
\end{equation}
Lei de Beer-Lambert (quantificação por FTIR):
\begin{equation}
    A = \varepsilon\,c\,l
\end{equation}
onde $A$ é a absorbância, $\varepsilon$ o coeficiente de absorção molar,
$c$ a concentração e $l$ o caminho óptico.

\subsubsection*{Importância}

FTIR e Raman são complementares: modos ativos em um tendem a ser inativos
no outro. Juntos, fornecem identificação completa de grupos funcionais,
análise de pureza, controle de qualidade industrial, estudo de
biomoléculas (proteínas, lipídios, ácidos nucleicos) e caracterização de
materiais como grafeno, nanotubos e polímeros.

% ---------------------------------------------------------------
\subsection{IRMS --- Espectrometria de Massas de Razão Isotópica}
% ---------------------------------------------------------------

\subsubsection*{Histórico}

A IRMS surgiu dos trabalhos de Nier e McKinney nos anos 1940--1950, que
construíram espectrômetros capazes de medir com precisão as razões entre
isótopos estáveis. A partir dos anos 1970, o acoplamento com analisadores
elementares (EA-IRMS) e cromatógrafos (GC-IRMS) expandiu enormemente as
aplicações, tornando a técnica essencial em geoquímica, arqueologia,
controle de alimentos e fisiologia.

\subsubsection*{Princípio de funcionamento}

A amostra é convertida em gases simples ($\text{CO}_2$, $\text{N}_2$,
$\text{SO}_2$, $\text{H}_2$) por combustão ou pirólise. Esses gases são
introduzidos no espectrômetro, ionizados por impacto de elétrons e
separados em um setor magnético. Coletores de Faraday simultâneos medem
as correntes iônicas dos isótopos de interesse com precisão da ordem de
0,01‰.

\subsubsection*{Grandezas e cálculos principais}

A razão isotópica da amostra é expressa em delta (\textperthousand)
relativa a um padrão internacional:
\begin{equation}
    \delta\,(\text{\textperthousand}) =
    \left(\frac{R_{\text{amostra}}}{R_{\text{padrão}}} - 1\right) \times 1000
\end{equation}
onde $R$ é a razão entre o isótopo pesado e o leve (ex.: $^{13}$C/$^{12}$C,
$^{15}$N/$^{14}$N, $^{18}$O/$^{16}$O).

\subsubsection*{Importância}

A IRMS permite rastrear a origem geográfica de alimentos e bebidas,
detectar adulterações, estudar ciclos biogeoquímicos, reconstruir paleoclimas
e investigar rotas metabólicas \textit{in vivo}. Na área de biocombustíveis,
é empregada para certificar a origem renovável de compostos e estimar
frações biogênicas em misturas com derivados fósseis.
